\section{Discussion}
\label{sec:discussion}
\subsection{Distinguishing supervision, forecast readout, and forecasting fidelity}
The seed-44 probability ensembles support a descriptive hierarchy in macro AUROC and AP: auxiliary supervision improves the direct-readout point estimates, and the complete forecast-attention readout improves them further. Only the complete-versus-NoJEPA macro AUROC interval excludes zero among these principal AUROC contrasts, but its nominal DeLong $p=0.0462$ becomes 0.4161 after Holm adjustment. None of the nine principal comparisons meets the adjusted criterion. Moreover, direct readout has higher macro F1 than forecast attention, and the mean-pooling control has the highest macro F1 among all four models. The ranking differences remain exploratory point-estimate advantages under this selected policy; they do not establish superiority or that attention is necessary.

The archived probes also challenge an explanation based solely on more accurate future states. At common epoch 20, the auxiliary-supervised direct models have better copy-normalized latent forecast skill than the full models, while the full probability ensemble has better pCR ranking. A possible explanation is that classification gradients shape forecast features toward endpoint-relevant directions at the expense of some fidelity to the EMA target. The nonzero pCR-to-predictor gradient is consistent with this pathway. However, fold-averaged training diagnostics in separately learned spaces cannot establish that mechanism or demonstrate anatomical forecasting quality. They are evaluated at a different level from the ensemble pCR predictor.

Useful future controls include classifier-width-matched direct readout, stopped-gradient forecast features, copy-state readout, and independent holdout forecast evaluation. The main models predict embeddings rather than generate MRI for classification. The internal name \nojepa{} denotes removal of the future-state objective while retaining shared fusion and transition queries. The forecast-versus-direct comparison also changes classifier width and feature groups, so it cannot isolate attention alone.

\subsection{What additional visits change}
The complete ensemble's AUROC increases across T0--T2 on the same patients, with a paired T2--T0 difference of +0.0099. Its F1 is nevertheless lower at T2 than at T0 and falls most at T1. The increased number of positive predictions identifies additional true positives but also adds false positives. These error counts explain why later information and improved ranking need not improve a fixed-threshold decision rule. Outcome-conditioned probability-change histograms describe ensemble score movement; they do not establish that every patient's prediction becomes more accurate.

T0 imaging preceded randomization, while the model is conditioned on the eventual assigned arm. This remains a retrospective conditional scenario. The arm encoding does not capture delivered dose, cycle timing, dose reductions, adherence, or treatment changes. Altering the arm token cannot identify an individual treatment effect without suitable causal assumptions and exposure data. Complete T0--T2 holdout eligibility also restricts the population to which the longitudinal comparison applies.

\subsection{Classification, calibration, and hypothetical utility}
The additional metrics reveal differences hidden by AUROC alone. Compared with direct readout, the complete model has higher point estimates of macro sensitivity and NPV but lower specificity, PPV, balanced accuracy, MCC, and F1. The sensitivity-difference interval includes zero, while the specificity-difference interval excludes zero conditionally on these fitted models. The latter is an unadjusted interval: its paired bootstrap test does not meet the secondary-family Holm criterion ($p_{\mathrm{Holm}}=0.14175$). All three principal models have macro raw accuracy below the all-negative reference of 0.68, although their AUROCs exceed 0.72. Raw accuracy would therefore obscure detection of the minority positive class.

The full model's NLL, Brier score, and ECE are also worse than the direct controls' point estimates. Probability ensembling does not by itself establish calibration. The OOF Youden analysis changes operating points, with small F1 increases for the principal models but a decrease for mean pooling. OOF development scores originate from one held-out fold model per patient, whereas holdout scores average five models. Transferring a threshold between these score distributions is a sensitivity analysis; it cannot replace independent ensemble calibration or validation of a clinical threshold.

Net-benefit rankings change with the hypothetical threshold. Direct readout has higher macro net benefit at 0.30, while forecast attention has higher net benefit at 0.70; all principal variants have negative macro values at 0.50. These are summaries under a generic positive-action utility, not evaluations of a defined care pathway. A clinically motivated operating point would require the consequences of false-positive and false-negative actions to be specified and calibration to be fitted without holdout leakage.

\subsection{Subgroup inference and transportability}
The four HR/HER2 groups have positive full-versus-control macro AUROC differences, but all paired intervals include zero. Their magnitudes differ and the smallest group has only 14 patients. These summaries provide no confirmed subgroup interaction or subtype-specific benefit. In HR-positive/HER2-negative patients, low sensitivity at 0.5 persists despite improved ranking; the group's low prevalence informs predictive values but does not itself explain sensitivity.

Source namespaces share one clinical release and do not supply independent external validation. Crop sensitivity analysis is similarly constrained: the 11 patients with partial clipping contain only one positive outcome, preventing reliable discrimination assessment under the stated minimum-outcome rule. Reporting the denominator and suppressing AUROC/AP is more informative than interpreting an unstable point estimate.

\subsection{Limitations and a confirmatory study}
\label{sec:limitations}
Seed 44 and probability ensembling were selected after holdout results were available. Patient bootstrap intervals condition on these fitted ensembles and fixed OOF thresholds. They exclude retraining, seed choice, aggregation choice, model selection, and threshold-selection uncertainty, and remain unadjusted for multiple comparisons. Separate Holm adjustments for the nine principal and 81 secondary tests, with an additional all-90 sensitivity analysis, do not remove retrospective selection bias. The testing methods and reporting families were also chosen retrospectively. DeLong inference is asymptotic, and null-centered bootstrap tests are approximate; finite-sample behavior can be problematic for non-smooth metrics such as AP, F1, and binned ECE. Neither adjusted family yielded a significant result, and nonsignificance does not demonstrate equivalence. One selected seed does not establish training stability. The reported predictor requires five fold models at inference, and its performance should not be interpreted as that of one model trained on the full development cohort. A confirmatory analysis should lock training, aggregation, and statistical policies before evaluating independent patients.

The holdout is an internal downstream split. Twenty-nine of its patients had belonged to an earlier development split. Current supervised models were retrained using the new patient-disjoint folds and were not warm-started from those earlier pCR models, but codec and localization-model pretraining have not been fully audited for current-holdout overlap. This prevents a claim that every pipeline stage was evaluated on wholly unseen patients.

Fixed baseline localization and cropping may omit multifocal disease or changes outside the field of view. Geometry and array checks do not establish anatomical registration accuracy. Holdout patients have complete T0--T2 imaging, whereas development retains incomplete histories; missingness, acquisition variation, and subtype prevalence may limit transportability. Archived training probes do not replace best-checkpoint holdout forecast evaluation and have small long-horizon denominators.

The study remains an internal architectural comparison. Matched clinical-only, conventional imaging, and independent longitudinal baselines are still needed, along with external validation, anatomical quality review, and independently assessed calibration. Precise pCR pathology criteria, scan intervals, ethics determination, and upstream provenance also require final documentation before submission.
